\section{Spectral Bundles}
\label{sec:spectral-bundles}

Spectral behavior under interrogation is modeled using a fiber bundle structure
over the cognitive manifold. This section introduces the spectral bundle
$\pi : \mathcal{B}_S \to \mathcal{M}_{cog}$, defines its fibers and sections,
and formalizes the geometric deformation induced by semantic excitation and
modal-boundary traversal.

\subsection{Bundle Structure}

The spectral bundle is defined as
\[
\pi : \mathcal{B}_S \to \mathcal{M}_{cog},
\]
where each fiber
\[
\mathcal{B}_{S,x} = \pi^{-1}(x)
\]
represents the spectral component space associated with interaction state
$x \in \mathcal{M}_{cog}$. The total space $\mathcal{B}_S$ encodes all possible
spectral emissions NDH may produce under interrogation.

\subsection{Spectral Fibers}

Each fiber $\mathcal{B}_{S,x}$ contains the spectral basis functions inherited
from NDH's runtime geometry. These basis functions form a structured set of
spectral directions whose alignment and deformation govern NDH's emission
behavior. The fiber structure ensures that spectral components vary smoothly
over $\mathcal{M}_{cog}$ except at modal boundaries.

\subsection{Sections and Emission}

A spectral section is a smooth map
\[
\Phi : \mathcal{M}_{cog} \to \mathcal{B}_S
\]
satisfying $\pi \circ \Phi = \mathrm{id}_{\mathcal{M}_{cog}}$. The value
$\Phi(x)$ represents the spectral emission associated with interaction state
$x$. Semantic excitation induces deformation of this section, captured by the
directional derivative
\[
D_{\mathbf{s}} \Phi_x,
\]
which describes the first-order spectral response to excitation vector
$\mathbf{s} \in T_x \mathcal{M}_{cog}$.

\subsection{Modal-Boundary Twisting}

Modal boundaries in $\mathcal{M}_{cog}$ induce twisting of the spectral bundle.
If $\gamma(t)$ is a geodesic crossing a modal boundary at $t = t_0$, then the
fiber transition map
\[
\tau_{\gamma(t_0)} : \mathcal{B}_{S,\gamma(t_0^-)} \to
                     \mathcal{B}_{S,\gamma(t_0^+)}
\]
may be nontrivial. Such twisting corresponds to structural changes in spectral
alignment, recursion influence, or non-dual algebraic coupling.

Modal-boundary twisting is a primary mechanism by which NDH exhibits
interaction-dependent spectral variation.

\subsection{Parallel Transport}

The interaction connection $\nabla^{cog}$ induces parallel transport on the
spectral bundle. Given a geodesic $\gamma(t)$, parallel transport defines a map
\[
P_\gamma : \mathcal{B}_{S,\gamma(0)} \to \mathcal{B}_{S,\gamma(T)},
\]
representing the evolution of spectral components along the trajectory. Parallel
transport interacts with modal-boundary twisting to produce nontrivial spectral
morphisms.

\subsection{Holonomy and Spectral Morphisms}

For closed interaction loops $\gamma$, parallel transport yields the spectral
holonomy operator
\[
\mathcal{H}_S(\gamma) : \mathcal{B}_{S,x} \to \mathcal{B}_{S,x},
\]
which captures the cumulative spectral deformation induced by traversing the
loop. Nontrivial spectral holonomy corresponds to interaction classes exhibiting
spectral drift, modal misalignment, or recursion-induced spectral rotation.

\subsection{Summary}

Spectral bundles provide the geometric structure governing NDH's spectral
response under interrogation. Semantic excitation deforms spectral sections,
modal boundaries twist fibers, and parallel transport induces spectral
morphisms. Holonomy of the spectral bundle yields a geometric classification of
interaction behavior, forming a central component of NDH interaction geometry.
